% ==============================================================================
% INTRODUCCIÓN
% ==============================================================================
% Instrucciones de la plantilla:
% - Proporcionar un contexto claro sobre los objetivos, fundamentos teóricos y relevancia de la práctica realizada.
% - Incluir explicación de los conceptos gráficos abordados en la práctica (ej: rasterización, transformaciones geométricas, algoritmos de dibujo, OpenGL/WebGL, shaders, etc.). Se recomienda usar la información investigada en los previos para elaborar esta sección.
% - ES OBLIGATORIO incluir citas a fuentes bibliográficas en formato APA.
%   Ejemplo: "De acuerdo con \cite{teodoro_vite}, el cómputo gráfico..."
% ==============================================================================

\section{Introducción}

El cómputo gráfico moderno se fundamenta en la manipulación y visualización eficiente de datos bidimensionales y tridimensionales mediante el cauce de renderizado (pipeline gráfico) acelerado por hardware \cite{hearn_baker}. En este contexto, librerías como OpenGL proporcionan una interfaz estándar y robusta para la comunicación directa con la unidad de procesamiento gráfico (GPU) \cite{khronos_opengl}.

De acuerdo con \cite{teodoro_vite}, el dominio de los fundamentos matemáticos y las técnicas de transformación geométrica resulta indispensable para la construcción de escenas interactivas y fotorrealistas. A lo largo de esta práctica, se exploran los principios de...

% Continuar la redacción integrando:
% 1. Contexto y motivación del tema.
% 2. Conceptos teóricos clave aplicados (rasterización, shaders, matrices MVP, etc.).
% 3. Citas bibliográficas correspondientes en formato APA.
